\documentclass[11pt,a4paper]{article}
\usepackage[T1]{fontenc}
\usepackage{lmodern,amsmath,amssymb,booktabs,microtype}
\usepackage[margin=24mm]{geometry}
\usepackage[colorlinks=true,urlcolor=blue,linkcolor=blue]{hyperref}
\newcommand{\EE}{\mathbb E}
\newcommand{\av}[1]{\langle #1\rangle}
\newcommand{\MM}{\widehat{\mathcal M}}
\newcommand{\Jh}{\mathcal J_h}
\DeclareMathOperator{\tr}{tr}
\title{Resolving the memory background by energy\\Conserved subtraction and the first fluctuation pair}
\author{Alan Day \and Chengheng\thanks{AI research assistant (OpenAI).}}
\date{Local research note --- 6 November 2026}
\begin{document}\maketitle
\begin{abstract}
The scalar mean of the temporal-singlet memory is insufficient as a
fluctuation counterterm. We determine its limiting regression against
energy functions using the established $q$-Mehler Markov operator.
An explicit positive variance gives a lower bound showing why scalar
centering fails at the collective normalization. Exact conserved
polynomials instead yield a finite pair of memory coefficients. Their
complete two-energy spectral matrix follows from the previously
computed quadratic, cubic and rank-filtered singlet channels. Energy
centering removes the time descendant present in the raw cubic
insertion; the normalized memory tangent tends to the cubic
$\operatorname{sech}^6$ kernel at small coupling.
\end{abstract}
\section{What the energy-weighted mean determines}
Use the same standard CAR complex $p$-body Gaussian Hamiltonian
$H=H_s$, including overlapping index sets, with
$w=4^p(J_s^2N/p^2)(p!/N^p)^2$. Brackets denote the annealed full
normalized trace at $\beta=\mu=0$, and $\hbar=1$. First take
$N,p\to\infty$, $p^2/N\to\lambda>0$, then optionally
$\lambda\downarrow0$. Set $q=e^{-4\lambda}$,
$\sigma^2=J_s^2/\lambda$; all times use string units.

The preceding stage defined, in each nonzero charge block,
\begin{equation}
 h=H_m,\quad K=1_N\otimes H_{m-1},\quad
 V\psi=m^{-1/2}\sum_i|i\rangle c_i\psi,\quad Q=1-VV^\dagger,
\end{equation}
and $\mathcal R=QKV$, $K_\perp=QKQ$. The energy-centered memory is
\begin{equation}
 \MM_m(s)=e^{-ihs/2}\mathcal R^\dagger e^{iK_\perp s}
 \mathcal R e^{-ihs/2},\qquad \MM_0(s)=0.
\end{equation}
Its scalar mean satisfies $p\av{\MM(s)}\to4J_s^2G_q(s)$.
For every bounded continuous energy test $f$, the stronger statement is
\begin{equation}\label{eq:regression}
 p\av{f(H)^\dagger\MM(s)}\longrightarrow
 4J_s^2\int\overline{f(x)}g_s(x)\,d\nu_q(x),\qquad
 g_s(x)=e^{-isx}(\mathsf P_q e^{is\cdot})(x).
\end{equation}
Here $\nu_q$ is the variance-$\sigma^2$ Hamiltonian law and
$\mathsf P_r$ is the established Mehler transition operator.

\newpage
\section{A nonconstant conserved background}
Write $d\mu_r(x,y)=d\nu_q(x)\,\mathsf P_r(x,dy)$, using the
same normalized two-energy probe measure as in the preceding stage.
For monic $q$-Hermite polynomials with
$H_{n+1}(x)=xH_n(x)-\sigma^2[n]_qH_{n-1}(x)$, the known Markov
operator obeys \cite{BKS}
\begin{equation}
 \mathsf P_rH_n=r^nH_n,\qquad
 \mathsf P_r^*=\mathsf P_r,\qquad \mathsf P_r\mathsf P_t=\mathsf P_{rt}.
\end{equation}
These are energy-space transition operators. Their parameter is not
the real memory time $s$. The Hamiltonian support is
$[-2\sigma/\sqrt{1-q},2\sigma/\sqrt{1-q}]$.

To prove \eqref{eq:regression}, insert $f(h)^\dagger$ into the two
memory--force comparison bounds from the preceding note. Their
right-hand sides acquire at most $\|f\|_\infty$. The charge-window
and tail estimates are unchanged. The resulting force expression is
an integral of $\overline{f(x)}e^{is(y-x)}$ against its positive
two-energy measure, which converges to $4J_s^2\mu_q$. This proves
the stated weighted mean. It does not prove strong operator
convergence to $4J_s^2g_s(H)$.

Nevertheless the nonconstant background is measurable exactly:
\begin{align}
 \int g_s\,d\nu_q&=G_q(s),\\
 V_{\rm bg}(s):=\int|g_s-G_q(s)|^2\,d\nu_q
 &=G_{q^2}(s)-G_q(s)^2\geq0.\label{eq:variance}
\end{align}
Indeed multiplication by $e^{-isx}$ preserves the norm, and
$\|\mathsf P_q e^{isx}\|_2^2
=\av{e^{isx},\mathsf P_{q^2}e^{isx}}_{L^2(\nu_q)}=G_{q^2}(s)$.
The kernels are real by simultaneous energy reflection. Low moments give
\begin{align}
 g_s(x)&=1+is(q-1)x
 -\tfrac12s^2\bigl[(1-q)^2x^2+(1-q^2)\sigma^2\bigr]+O(s^3),\\
 V_{\rm bg}(s)&=\sigma^2(1-q)^2s^2+O(s^4).
\end{align}
Thus it is strictly positive for sufficiently small nonzero $s$.

Let $X_N=p\MM(s)$ and $\|B\|_2^2=\av{B^\dagger B}$. Taking the
test $g_s-G_q(s)$ in \eqref{eq:regression}, with a bounded continuous
extension outside the limiting support, and applying Cauchy--Schwarz,
\begin{equation}
 \liminf_{\rm DS}\|X_N-\av{X_N}1\|_2^2
 \geq16J_s^4V_{\rm bg}(s).\label{eq:lower}
\end{equation}
The full-trace charge number is $Q_f$. Central-window concentration
and $\|\MM\|_2^2\leq\av{G^2}=O(p^{-2})$ give
\begin{equation}\label{eq:divergence}
 \liminf_{\rm DS}\frac{\|Q_f\MM-\av{Q_f\MM}1\|_2^2}{p^2}
 \geq\frac{4J_s^4}{\lambda^2}V_{\rm bg}(s).
\end{equation}
Scalar centering therefore leaves a divergent collective background.

\newpage
\section{Exact subtraction at the first two memory orders}
At finite $N$, set $r=p/m$ and expand
\begin{equation}
 \MM_m(s)=G+isZ_1+O(s^2),\qquad
 G=A_m/m-r^2h^2,\qquad
 mZ_1=T_{3,m}+2r\Jh A_m-mr^3h^3.\label{eq:jet}
\end{equation}
Here $A=\sum_i[H,c_i]^\dagger[H,c_i]$,
$M_3=\sum_i c_i^\dagger\operatorname{ad}_H^3c_i$,
$T_3=(M_3+M_3^\dagger)/2$, and $\Jh B=(hB+Bh)/2$.
Let $C_m=\av{h^2}_m$, $a_m=\av{A_m}_m$, and use the previously
defined ensemble projection coefficient
\begin{equation}
 c_m=\frac{\av{(A_m-a_m)(h^2-C_m)}_m}
 {\av{(h^2-C_m)^2}_m},\qquad
 Y_m=A_m-a_m-c_m(h^2-C_m).
\end{equation}
All scalar coefficients here are deterministic charge-block functions.
For the already calculated cubic channels put
\begin{align}
 a&=-\frac{(1-q)^2}{1+q^2},&
 b&=\frac{2\lambda(1-q)(1+q^3)}{(1+q^2)(1+q+q^2)},\\
 R&=W_3T_3-aW_3(\Jh A)-bW_3h^3,\\
 F&=W_1T_3-\bar t_mh,&
 B_S&=W_1(\Jh A)-\bar s_mh,\quad B_Z=W_1h^3-\bar z_mh.
\end{align}
$W_j$ denotes ordinary Gaussian coupling chaos; bars are projections
of the respective linear-chaos operator onto $h$. The earlier exact
cubic reduction gives
\begin{align}
 P_m&=(ac_m+b)h^3+[\bar t_m-(ac_m+b)\bar z_m]h,\\
 T_{3,m}-P_m&=R+a\Jh Y_m+F-aB_S-bB_Z.\label{eq:cubic}
\end{align}
The required two conserved counterterms are consequently explicit:
\begin{align}
 C_{0,m}&=a_m+c_m(h^2-C_m)-mr^2h^2,\\
 C_{1,m}&=P_m+2r[a_mh+c_m(h^3-C_mh)]-mr^3h^3.\label{eq:C1}
\end{align}
Both commute with $h$. Subtracting them gives the exact pair
\begin{align}
 X_{0,N}&=mG-C_{0,m}=Y_m,\\
 X_{1,N}&=mZ_1-C_{1,m}
 =R+(a+2r)\Jh Y_m+F-aB_S-bB_Z.\label{eq:X1}
\end{align}
At fixed $\lambda$, the previous rank-tail bounds imply
$B_S,B_Z=O_{L^2}(p^{-1})$ and $r\Jh Y\to0$. Hence the second
coefficient has limiting representative $L=R+a\mathcal J_HY+F$.
This names a limit of moments and two-point spectra, not an identification
of Hilbert spaces at different $N$.

\newpage
\section{The full two-energy spectral matrix of the pair}
Use the established vanishing cross measures for $F,Y,R$ and their
diagonal measures $V_F\mu_q$, $W\mu_{q^2}$, $V_R\mu_{q^3}$.
For $e=(x+y)/2$ the pair $(Y,L)$ has the positive spectral matrix
\begin{equation}\label{eq:matrix}
 d\boldsymbol\nu=
 \begin{pmatrix}
 W\,d\mu_{q^2}&a eW\,d\mu_{q^2}\\
 a eW\,d\mu_{q^2}&V_F\,d\mu_q+V_R\,d\mu_{q^3}
                  +a^2e^2W\,d\mu_{q^2}
 \end{pmatrix}.
\end{equation}
Relative to the product Hamiltonian measure, its determinant equals
$WK_{q^2}(V_FK_q+V_RK_{q^3})\geq0$. The coefficients are
\begin{align}
 W&=4J_s^4\left(\lambda^{-1}+\frac{4q}{1+q}\right),&
 V_F&=64J_s^6q^2,\\
 V_R&=\sigma^6\frac{\lambda A_1+\lambda^2A_2}{D^2},&
 D&=(1+q^2)(1+q+q^2),\\
 A_1&=4(1-q^2)(1+q^2)(1+q+q^2)^2,\nonumber\\
 A_2&=16q^2(1+q)(1+q^2)(1+q+q^2)(2q^2-q+2).\nonumber
\end{align}
These diagonal channel results are inherited from the preceding
singlet calculations, not recalculated as new probe formulas.
The application here is the exact memory-coefficient combination.

At any external centre time $T$, the off-diagonal ordinary correlator
vanishes. The transformation $(x,y)\mapsto(-y,-x)$ preserves $y-x$
and the measure but reverses $e$. Therefore
\begin{align}
 \av{Y^\dagger L(T)}_{\rm DS}&=0,\\
 \av{Y^\dagger Y(T)}_{\rm DS}&=WG_{q^2}(T),\\
 \av{L^\dagger L(T)}_{\rm DS}&=V_FG_q(T)+V_RG_{q^3}(T)+a^2WJ_2(T),\\
 J_2(T)&=\int e^2e^{iT(y-x)}\,d\mu_{q^2}(x,y).
\end{align}
Ordinary-time orthogonality does not erase the nonzero energy-resolved
cross entry in \eqref{eq:matrix}. Its energy-weighted integral is
$aW\sigma^2(1+q^2)/2$.

The first-degree memory polynomial $\Xi(s)=Y+isL$ now has the
complete centre-time covariance
\begin{equation}
 \av{\Xi(s)^\dagger\Xi(u;T)}=
 WG_{q^2}(T)+su\,[V_FG_q(T)+V_RG_{q^3}(T)+a^2WJ_2(T)].
\end{equation}
This is an exact covariance of the limiting pair. For the actual
finite-memory field it specifies the $(0,0),(1,0),(0,1),(1,1)$
Taylor coefficients, with derivatives taken before double scaling.
Terms of degree two in either memory argument remain uncomputed.

\newpage
\section{What energy centering selects at small coupling}
The original non-Hermitian derivative has
\begin{equation}
 M_3-P_m\longrightarrow R+a\mathcal J_HY+F+\tfrac32[H,Y].
\end{equation}
The memory tangent $L$ has no final time-descendant term. This follows
from the exact centered memory expansion, not from manually discarding
a large contribution to the raw cubic spectrum.

Its squared norm and second frequency moment are
\begin{align}
 N_L&=V_R+V_F+a^2W\sigma^2(1+q^2)/2,\\
 D_L&=2\sigma^2\{V_R(1-q^3)+V_F(1-q)\}
       +\tfrac12 a^2W\sigma^4(1+q)(1-q^4).
\end{align}
In the subsequent $\lambda\downarrow0$ limit,
\begin{equation}
 \lambda W\to4J_s^4,\qquad \lambda N_L\to32J_s^6,
 \qquad D_L/N_L\to24J_s^2.
\end{equation}
The $F$ fraction is $2\lambda+o(\lambda)$; the energy-descendant
fraction is $8\lambda^3+o(\lambda^3)$. Thus
\begin{equation}
 \frac{\av{L^\dagger L(T)}}{N_L}\longrightarrow
 \operatorname{sech}^6(2J_sT),\qquad
 \frac{\av{Y^\dagger Y(T)}}W\longrightarrow
 \operatorname{sech}^4(2J_sT).
\end{equation}
This applies the established small-coupling probe kernels to the
specific memory coefficients. The two labels do not by themselves
determine bulk meson masses.

\section{Verification and actual finite-coupling effects}
Three finite CAR systems verify five exact rational-complex matrix
identities each: the two subtractions, conservation of both
counterterms, and Hermiticity of the tangent. The calculation obtains
the first-return coefficient directly from $V^\dagger K^nV$ and obtains
the cubic chaos components separately by Gaussian polarization.
Twenty-eight exact polynomial tests verify the conditional
$q$-Hermite action through total degree six; symbolic Laurent limits
give $4,32,768$ for $J_s=1$ and the three unnormalized limits above.

For four positive couplings we evaluate the background variance and
tangent kernel by Jacobi evolution, with a larger truncation as a
check. At $\lambda=0.25,0.1$, independent 128/192-node Mehler
integration also checks the conditional variance directly, the zero
ordinary cross and nonzero energy-weighted cross, and the tangent
kernel. Finite-$\lambda$ differences from $\operatorname{sech}^6$
are retained; they are not numerical errors.

One initial independent comparison failed: the quadrature code had
omitted $K_{q^2}$ on the $e^2$ descendant term, giving a nonunit
normalized tangent already at $T=0$. The original source and values
are preserved. Restoring that factor fixes the formula; no tolerance
was relaxed. Floating-point refinement is not interval certification.

\newpage
\section{From the tangent pair to finite memory}
The next unknown is the full covariance after conserved subtraction,
with two finite memory parameters and an independent centre time.
The present work determines its first pair of coefficients and proves
that a purely scalar subtraction cannot yield the desired scaling.
Equation~\eqref{eq:regression} also gives the full leading
energy-dependent counterterm as a regression functional, at every
fixed memory time.

Three requirements now have a concrete form. First, compute the
next centered return coefficient, including the $G^2$ return term
and the fourth derivative singlet, then determine its conserved
polynomial and its mixed spectra with $Y,L$. Second, control the
operator remainder with two memory insertions at the collective
normalization. The single-insertion comparison used for the mean
does not provide this bound. Third, match the resulting microscopic
kernel with the bulk internal collective operator and spatial
insertions, preserving the distinction between memory time, centre
time and a light-front quark fraction.

The energy-dependent regression $g_s$ alone is not the full finite-$N$
conditional expectation onto the commutant. Nor does subtracting its
leading $1/p$ term automatically fix all subleading conserved pieces
at the $m\sim p^2$ normalization. Equations~\eqref{eq:C1}--\eqref{eq:X1}
show explicitly how such pieces are fixed at the first two orders.
At those orders the limiting spectra are absolutely continuous and
have no zero-frequency atom. Finite-memory tightness remains an
additional problem.

The complete fluctuation field and the spatial holographic dictionary
remain open. The useful advance is a specified conserved background,
an obstruction to scalar centering, and a positive spectral matrix
for the actual first two memory coefficients.

\paragraph{Local record.}
Source, exact checks, numerical parameters, dependency versions,
initial failure and input hashes are retained in the project stage
directory. The PDF is compiled twice and reviewed page by page;
its source/output hashes are bound to the local evidence record.
No Git commit, push or release archive is created in this stage.

\begin{thebibliography}{9}
\bibitem{BKS} M.~Bo\.{z}ejko, B.~K\"ummerer and R.~Speicher,
\emph{$q$-Gaussian processes: non-commutative and classical aspects},
\href{https://arxiv.org/abs/funct-an/9604010}{arXiv:funct-an/9604010},
Theorems~1.10 and 4.6, especially pp.~20--22.
\bibitem{BNR} M.~Berkooz, V.~Narovlansky and H.~Raj,
\emph{Complex Sachdev--Ye--Kitaev model in the double scaling limit},
\href{https://arxiv.org/abs/2006.13983v2}{arXiv:2006.13983v2}.
\bibitem{local} A.~Day and Chengheng, local project notes:
\emph{From the charged force to the mean memory kernel};
\emph{Conserved-subtracted cubic spectrum}; and the associated
singlet-kernel, cubic-cross and filtered-spectrum calculations,
6 November 2026. These supply the model-specific precursor identities
and measures used here; no priority claim is made.
\end{thebibliography}
\end{document}
